% 2026年8月月度报告；XeLaTeX 编译两遍。
\documentclass[stage1,fontset=fandol]{ydzjreport}
\stageonesetup
  {阶段一 · 文献复现与结构探索}
  {冯旭东}
  {流道形态}
  {2026-08-01 至 2026-08-31}
  {文献复现；分区冷却；流道形态；综合性能评价}
  {复现Ao等的三种微通道并获得一致的变化趋势，借鉴Yin的评价方法与Seong的结构思路，整理23项候选结构记录。}
\begin{document}
\makestageonehead

\section{阅读范围}
\subsection{文献清单}
本月围绕微通道结构设计、流动与换热性能评价开展文献阅读，并结合项目申报书及前期仿真记录推进结构探索。本人作为流道形态制作的主要成员，重点承担几何构造、仿真复现和候选方案比较。主要阅读文献见表~\ref{tab:reading}。
\begin{table}[H]
\centering
\caption{本阶段主要阅读文献及工作关联}
\label{tab:reading}
\begin{tabularx}{\textwidth}{@{}c Y Y@{}}
\toprule
\rowcolor{ydzjlight} 序号 & 文献与研究对象 & 对本阶段工作的作用 \\
\midrule
1 & Ao等：方形、圆形与六边形微通道散热器\cite{ao} & 复现三种流道，学习短流路、多进出口与局部强化的设计方法 \\
2 & Yin等：同轴开环环形TSV针肋微通道\cite{yin} & 借鉴数据分析方法，建立换热与阻力共同评价的口径 \\
3 & Seong等：直列与交错布置的倾斜肋--针肋微通道\cite{seong} & 获得倾斜、交错及肋--针肋组合等结构探索思路 \\
\bottomrule
\end{tabularx}
\end{table}

\subsection{选读理由}
项目面向空间与时间上不均匀的热载荷，需要为不同区域配置适合的冷却结构。Ao等的研究将流路布局与热点区强化相结合，适合作为几何建模和仿真复现的起点。Yin等将最高温度、平均努塞尔数、压降及综合性能指标结合分析，帮助判断哪些数据能够反映流道实际表现。Seong等从肋倾角和排列方式出发研究流动重分配，为前期多个形态的设计提供了新的思路。三篇文献分别支撑结构复现、性能评价和方案拓展。

\clearpage
\section{核心概念梳理}
\subsection{概念一：}
\textbf{非均匀热载荷与分区适配。} 局部发热强度、温度水平和温度余量不同，意味着冷却资源应按区域配置。高温热点需要重点控制峰值温度，低负荷区域需要减少不必要的流动阻力，变化区则需要兼顾区域过渡与负荷变化。

分小块的目的，是将空间不均匀的热需求转化为可单独研究和组合的流道单元。各小块可在连续冷板上形成不同几何区域，并不必然意味着将冷板机械切开或完全隔热。先研究单元、再研究组合，有利于识别各形态的作用，并为独立供液或流量调控预留条件。

\subsection{概念二：}
\textbf{能真实反映流道表现的数据与评价口径。} 借鉴Yin等的分析\cite{yin}，本阶段认识到不能仅用某一个温度值或出口温升判断流道优劣：固体最高温度$T_{\max}$反映热点约束，指定表面的平均温度与温差反映热状态及均匀性，平均努塞尔数$Nu_{\mathrm{avg}}$反映对流换热，进出口压差$\Delta p$反映水力代价。

Yin等采用的综合性能评价准则为
\[
 \mathrm{PEC}=\frac{Nu/Nu_0}{(f/f_0)^{1/3}},
\]
其中$Nu$和$f$分别为努塞尔数与摩擦因子，下标0表示基准直微通道。使用该指标需统一基准、雷诺数、特征尺度和物性口径，不能将不同条件下的温降与压差直接拼成PEC。水力功率还可按
\[
 P_{\mathrm h}=\Delta p\,Q_{\mathrm v}
\]
估算，其中$Q_{\mathrm v}$为体积流量；该值不包含泵效率及外部管路损失。后处理需说明取值位置、平均方式及保存时刻。固体温度、加热面温度和出口流体温度具有不同含义；面积平均与体积平均也不能互换。温度均匀性可在同一指定表面用$T_{\max}-T_{\min}$或温度标准差表征，不能用最高温度减任意平均温度替代完整评价。

\section{阶段一：通读与摘录}
\subsection{学习路径}
\begin{enumerate}
\item 通读Ao等论文，梳理三种流道的几何布局、进出口组织、边界条件和评价指标，完成三种结构的阅读与仿真复现。
\item 结合Yin等的分析，明确温度、压降、平均努塞尔数和PEC的用途，将评价重点由单纯降温转向换热与阻力的共同比较。
\item 阅读Seong等关于倾斜肋--针肋、底壁直列及上下壁交错布置的研究，关注流体混合、边界层扰动与流动停滞区的变化。
\item 将文献启发用于前期结构探索，整理23项候选结构或编号方案，其中20项已有最高温度、15项已有压差记录，形成后续筛选依据。
\end{enumerate}

\subsection{关键摘录}
表~\ref{tab:notes}对三篇论文要点作概括性转述，并列出与本人工作的关联。
\begin{table}[H]
\centering
\caption{文献要点与结构研究启发}
\label{tab:notes}
\begin{tabularx}{\textwidth}{@{}c Y Y@{}}
\toprule
\rowcolor{ydzjlight} 文献 & 论文要点 & 本人的理解与应用 \\
\midrule
Ao等\cite{ao} & 三种短流路、多进出口布局结合针肋扰动，比较结温、压降和温度均匀性 & 从整体流路和局部强化两个层面构造形态，并通过复现核对性能趋势 \\
Yin等\cite{yin} & 用平均努塞尔数、最高温度、压降和PEC比较不同针肋结构及排列 & 温度优势必须与流动代价共同评价；先明确数据含义，再进行结构排序 \\
Seong等\cite{seong} & 肋倾角及上下壁交错排列影响纵向流动重分配与换热--阻力平衡 & 除改变截面外，还可调整倾斜角、肋高与排列方式，拓展复合形态 \\
\bottomrule
\end{tabularx}
\end{table}

\section{阶段二：导出与复现}
\subsection{复现目标}
本阶段已阅读并复现Ao等论文中的\textbf{方形、圆形和六边形三种微通道}，取得与论文一致的结果变化趋势。复现旨在掌握三种布局的建模与流动换热分析方法，并为自主形态设计建立参照。

同时整理前期《仿真条件与数据》中的多形态探索记录，进行温度单位换算与压差复核。Ao三种结构的文献复现和前期自主方案的数据整理是本月相互衔接的两项工作。

\subsection{复现环境与步骤}
三种流道采用CFD方法开展流动与换热数值复现，参考Ao等论文的物理模型及数值方法\cite{ao}；前期数据整理采用Python和openpyxl。具体工作如下：
\begin{enumerate}
\item 根据论文的方形、圆形和六边形布局构造流道，梳理流体域、固体域、加热区域及进出口位置。
\item 对照论文的边界条件、物性和流动换热求解方法开展计算，检查残差以及温度、压差等监测量，再整理可用于比较的结果。
\item 对三种结构的结果进行趋势对照，关注热负荷及流量变化下的温度响应，以及强化换热伴随的流动代价。
\item 对前期数据按$T_{\mathrm C}=T_{\mathrm K}-273.15$换算温度，按$\Delta p=p_{\mathrm{in}}-p_{\mathrm{out}}$复算压差；保留网格风险及求解备注，对缺失结果不填零。
\end{enumerate}

\subsection{复现结果}
\textbf{Ao三种流道的仿真复现已完成，计算结果与论文趋势一致。} 三种结构及复现工作见表~\ref{tab:repro}。论文显示，在比较范围内，质量流量增加时结温总体下降、压降增加；热流密度增加时结温上升；方形结构在换热与阻力的综合比较中表现较好\cite{ao}。本阶段的复现支持以这些趋势作为后续结构研究的参照。当前报告记录趋势一致性，逐点数值与误差统计待随复现数据表归档。
\begin{table}[H]
\centering
\caption{Ao论文三种流道的复现工作}
\label{tab:repro}
\begin{tabularx}{\textwidth}{@{}Y Y Y@{}}
\toprule
\rowcolor{ydzjlight} 流道形态 & 复现关注内容 & 本阶段结果 \\
\midrule
方形微通道（SMHS） & 方形流路及中心扰动结构对流动与换热的影响 & 已完成复现，结果变化趋势与论文一致 \\
圆形微通道（CMHS） & 环形流路及流体分配对温度和压差的影响 & 已完成复现，结果变化趋势与论文一致 \\
六边形微通道（HMHS） & 六边形布局及局部流动组织的作用 & 已完成复现，结果变化趋势与论文一致 \\
\bottomrule
\end{tabularx}
\end{table}

在此基础上，前期探索覆盖平直、多条带、正弦、圆柱与六边形针肋、高低横向肋、三角肋、倾斜与交错针肋等。Seong等提供了通过倾斜、交错与肋--针肋组合促进混合的设计启发；上述多个形态作为自主候选方案记录，不将其全部等同于Seong论文的原始模型。代表性初步结果见表~\ref{tab:earlyperf}。

\begin{table}[H]
\centering
\caption{前期自主探索结构的温度与压差记录}
\label{tab:earlyperf}
\begin{tabularx}{\textwidth}{@{}Y c c Y@{}}
\toprule
\rowcolor{ydzjlight} 结构 & 最高温度 / $^\circ$C & 压差 / kPa & 筛选认识 \\
\midrule
Plain基准 & 72.02 & 1.169 & 流阻较小，峰值温度偏高 \\
Straight channel\_2 & 78.61 & 13.737 & 不能凭直线外形预判性能 \\
高低横向肋 & 53.08 & 58.824 & 控温改善伴随压降增加 \\
三角高低横向肋 & 39.86 & 662.360 & 阻力很高，网格标记较危险 \\
\bottomrule
\end{tabularx}
\end{table}

这些初步记录体现出降温与流阻之间的取舍，但各记录的边界条件尚未完全统一，不能据此给出形态的独立改进幅度。后续将借鉴Yin等的评价口径，在可比工况下重新筛选。

\section{阶段性收获}
\subsection{认知上的变化}
本月完成了从文献阅读、三种结构复现到自主方案探索的衔接。Ao等帮助建立几何布局与性能变化之间的认识，Yin等使数据分析从单看温度转向换热、压降及综合指标的共同判断，Seong等则拓展了倾斜、交错和复合结构的设计思路。

已形成23项候选结构记录及低温区、变化区、高温区的分区研究框架。下一步由多形态探索收敛到少数具有代表性的结构，结合温度余量与水力代价明确各自适用的冷却区域。

\subsection{尚存的疑问}
\begin{enumerate}
\item Ao三种结构的逐点复现误差及网格敏感性有多大？需将对应工况、原文取值和计算数据统一归档，补齐定量对照。
\item 前期网格危险提示及温度振荡是否影响结果排序？需检查局部网格与监测历史，并统一流动模型及边界条件。
\item 不同形态如何统一取值口径？原表PEC与EER的公式及基准需重新核对。
\item 变化区应如何识别？空间过渡带应关注温度梯度，动态负荷区应关注温度随时间的变化，两者需分开验证。
\end{enumerate}

\section{下月计划}
2026年9月拟将多结构探索收敛到少数代表性形态，具体计划见表~\ref{tab:plan}。
\begin{table}[H]
\centering
\caption{下月工作计划}
\label{tab:plan}
\begin{tabularx}{\textwidth}{@{}Y Y c@{}}
\toprule
\rowcolor{ydzjlight} 事项 & 预期产出 & 时间 \\
\midrule
收敛到直线、正弦、NACA2416和gao2.1四种形态 & 模型与工况清单 & 第1周 \\
整理不同热流密度下的温度和压差，核对单位与取值口径 & 可追溯数据表及比较曲线 & 第2--3周 \\
结合温度余量与流动代价提出候选分区方案 & 低温区、变化区和高温区配置建议 & 第4周 \\
\bottomrule
\end{tabularx}
\end{table}

\begin{thebibliography}{9}
\bibitem{ao} Ao C, Xu B, Chen Z. Design and performance optimization of novel microchannel heat sinks for high-power electronic chips: A comparative study on geometric configurations and thermal management efficiency. International Communications in Heat and Mass Transfer, 2026, 178: 111803. \href{https://doi.org/10.1016/j.icheatmasstransfer.2026.111803}{DOI}.
\bibitem{yin} Yin Y, Wang Z, Ma S, et al. A novel coaxial open annular TSV pin-fin embedded in microchannels to comprehensively improve the thermal stability, signal integrity, and heat transfer of 3D integrated circuits. International Communications in Heat and Mass Transfer, 2026, 172: 110619. \href{https://doi.org/10.1016/j.icheatmasstransfer.2026.110619}{DOI}.
\bibitem{seong} Seong H B, Gwon J G, Lee M S, et al. Numerical study and ANN analysis on heat transfer and flow characteristics of microchannel heat sink with inclined rib--pin fin in in-line and staggered arrangements. International Communications in Heat and Mass Transfer, 2026, 178: 111593. \href{https://doi.org/10.1016/j.icheatmasstransfer.2026.111593}{DOI}.
\end{thebibliography}
\end{document}
